% GENERATED FILE. Edit canonical lessons, then run npm run generate:books.
% canonical-source-hash: fnv1a64:d8b6d6c8b53ec83b
% canonical-lessons: ZH-C305-tanlianai, ZH-C305-qiuhun, ZH-C305-dinghun, ZH-C305-xishi, ZH-C305-xinlang

\chapter{Getting Engaged}
\label{ch:zh-getting-engaged}
\hlchaptermodality{305}

\begin{chapteropening}
\textbf{By the end of this chapter:} \emph{I can talk about going out with someone, a proposal, an engagement and a happy family occasion.}

Complete the last lesson of chapter 305: I can talk about going out with someone, a proposal, an engagement and a happy family occasion.
\end{chapteropening}

\section[tán liàn'ài]{tán liàn'ài — to go out with someone; to be in a relationship}

\label{lesson:ZH-C305-tanlianai}



\begin{quote}
Before the new one: say the Mandarin for school tuition fees, then the Mandarin for to commit to memory.
\end{quote}

\subsection*{You'll want to know: tán liàn'ài}
\emph{tán liàn'ài} — \textquotedblleft{}to go out with someone; to be in a relationship\textquotedblright{}.

\begin{quote}
\emph{tāmen zài tán liàn'ài} — They're going out together.
\end{quote}

Its characters come later; for now it is a word for the ear.

\begin{cousinweb}
\emph{tán liàn'ài} is built from \emph{tán} 'to talk', \emph{liàn} 'to long for' and \emph{ài} 'love', all inherited from Middle Chinese, \emph{dam}, \emph{ljwenH} and \emph{'ojH} (Guangyun, in Baxter's transcription), which glosses the first 'to talk', the second 'to long for' and the third 'to cherish': to talk love.
\end{cousinweb}

\begin{grammarlens}[title={What you've built}]
One more word for telling what happened.
\end{grammarlens}

\subsection*{Your turn}
\begin{itemize}
\raggedright
  \item \emph{Say it:} \emph{tán liàn'ài}
  \item \emph{Say it:} \emph{tán liàn'ài}, once more
  \item \emph{Say it:} \emph{tāmen zài tán liàn'ài} — They're going out together.
  \item \emph{Recall:} say the Mandarin for school tuition fees, then the Mandarin for to commit to memory, then say \emph{tán liàn'ài} again
\end{itemize}

\begin{tcolorbox}[breakable,colback=teal!4,colframe=teal!35!black,title={Before you move on}]
What does tán liàn'ài mean? (\textquotedblleft{}To go out with someone; to be in a relationship\textquotedblright{}.) Say it once more.
\end{tcolorbox}
\section[qiúhūn]{qiúhūn — to propose marriage}

\label{lesson:ZH-C305-qiuhun}



\begin{quote}
Before the new one: say the Mandarin for to commit to memory, then the Mandarin for to go out with someone.
\end{quote}

\subsection*{You'll want to know: qiúhūn}
\emph{qiúhūn} — \textquotedblleft{}to propose marriage\textquotedblright{}.

\begin{quote}
\emph{tā zuótiān wǎnshang qiúhūn le} — He proposed last night.
\end{quote}

Its characters come later; for now it is a word for the ear.

\begin{cousinweb}
\emph{qiúhūn} is built from \emph{qiú} 'to seek' and \emph{hūn} 'marriage', both inherited from Middle Chinese, \emph{gjuw} and \emph{xwon} (Guangyun, in Baxter's transcription), which glosses the first 'to search for' and the second 'marriage'.
\end{cousinweb}

\begin{grammarlens}[title={What you've built}]
One more word for telling what happened.
\end{grammarlens}

\subsection*{Your turn}
\begin{itemize}
\raggedright
  \item \emph{Say it:} \emph{qiúhūn}
  \item \emph{Say it:} \emph{qiúhūn}, once more
  \item \emph{Say it:} \emph{tā zuótiān wǎnshang qiúhūn le} — He proposed last night.
  \item \emph{Recall:} say the Mandarin for to commit to memory, then the Mandarin for to go out with someone, then say \emph{qiúhūn} again
\end{itemize}

\begin{tcolorbox}[breakable,colback=teal!4,colframe=teal!35!black,title={Before you move on}]
What does qiúhūn mean? (\textquotedblleft{}To propose marriage\textquotedblright{}.) Say it once more.
\end{tcolorbox}
\section[dìnghūn]{dìnghūn — to get engaged}

\label{lesson:ZH-C305-dinghun}



\begin{quote}
Before the new one: say the Mandarin for to go out with someone, then the Mandarin for to propose marriage.
\end{quote}

\subsection*{You'll want to know: dìnghūn}
\emph{dìnghūn} — \textquotedblleft{}to get engaged\textquotedblright{}.

\begin{quote}
\emph{tāmen shàng ge yuè dìnghūn le} — They got engaged last month.
\end{quote}

Its characters come later; for now it is a word for the ear.

\begin{cousinweb}
\emph{dìnghūn} is built from \emph{dìng} 'to agree on' and \emph{hūn} 'marriage', both inherited from Middle Chinese, \emph{dengX} and \emph{xwon} (Guangyun, in Baxter's transcription), which glosses the first 'to discuss and settle' and the second 'marriage'.
\end{cousinweb}

\begin{grammarlens}[title={What you've built}]
One more word for telling what happened.
\end{grammarlens}

\subsection*{Your turn}
\begin{itemize}
\raggedright
  \item \emph{Say it:} \emph{dìnghūn}
  \item \emph{Say it:} \emph{dìnghūn}, once more
  \item \emph{Say it:} \emph{tāmen shàng ge yuè dìnghūn le} — They got engaged last month.
  \item \emph{Recall:} say the Mandarin for to go out with someone, then the Mandarin for to propose marriage, then say \emph{dìnghūn} again
\end{itemize}

\begin{tcolorbox}[breakable,colback=teal!4,colframe=teal!35!black,title={Before you move on}]
What does dìnghūn mean? (\textquotedblleft{}To get engaged\textquotedblright{}.) Say it once more.
\end{tcolorbox}
\section[\texorpdfstring{\zh{喜事} (xǐshì)}{喜事 (xǐshì)}]{\zh{喜事} (xǐshì) — a happy occasion; a wedding in the family}

\label{lesson:ZH-C305-xishi}



\begin{quote}
Before the new one: say the Mandarin for to propose marriage, then the Mandarin for to get engaged.
\end{quote}

\subsection*{You'll want to know: \zh{喜事}}
\textbf{\zh{喜事}} — \emph{xǐshì} — \textquotedblleft{}a happy occasion; a wedding in the family\textquotedblright{}.

\begin{quote}
\textbf{\zh{我们家有喜事了}} \emph{wǒmen jiā yǒu xǐshì le} — There's a happy occasion in our family.
\end{quote}

\begin{cousinweb}
\textbf{\zh{喜事}} (\emph{xǐshì}) is built from \zh{喜} 'joy' and \zh{事} 'a matter', both inherited from Middle Chinese, \emph{xiX} and \emph{dzriH} (Guangyun, in Baxter's transcription), which glosses the first 'joy'.
\end{cousinweb}

\begin{grammarlens}[title={What you've built}]
One more word for telling what happened.
\end{grammarlens}

\subsection*{Your turn}
\begin{itemize}
\raggedright
  \item \emph{Say it:} \emph{xǐshì}
  \item \emph{Say it:} \emph{xǐshì}, once more
  \item \emph{Say it:} \emph{wǒmen jiā yǒu xǐshì le} — There's a happy occasion in our family.
  \item \emph{Recall:} say the Mandarin for to propose marriage, then the Mandarin for to get engaged, then say \emph{xǐshì} again
\end{itemize}

\begin{tcolorbox}[breakable,colback=teal!4,colframe=teal!35!black,title={Before you move on}]
What does \zh{喜事} mean? (\textquotedblleft{}A happy occasion; a wedding in the family\textquotedblright{}.) Say it once more.
\end{tcolorbox}
\section[xīnláng]{xīnláng — bridegroom}

\label{lesson:ZH-C305-xinlang}



\begin{quote}
Before the new one: say the Mandarin for to get engaged, then the Mandarin for a happy occasion.
\end{quote}

\subsection*{You'll want to know: xīnláng}
\emph{xīnláng} — \textquotedblleft{}bridegroom\textquotedblright{}.

\begin{quote}
\emph{xīnláng jīntiān hěn shuài} — The groom is very handsome today.
\end{quote}

Its characters come later; for now it is a word for the ear.

\begin{cousinweb}
\emph{xīnláng} is built from \zh{新} 'new' and \emph{láng}, both inherited from Middle Chinese, \emph{sin} and \emph{lang} (Guangyun, in Baxter's transcription), which glosses the second as the title of an official; its use for a young man is not given there.
\end{cousinweb}

\begin{grammarlens}[title={What you've built}]
One more word for telling what happened.
\end{grammarlens}

\subsection*{Your turn}
\begin{itemize}
\raggedright
  \item \emph{Say it:} \emph{xīnláng}
  \item \emph{Say it:} \emph{xīnláng}, once more
  \item \emph{Say it:} \emph{xīnláng jīntiān hěn shuài} — The groom is very handsome today.
  \item \emph{Recall:} say the Mandarin for to get engaged, then the Mandarin for a happy occasion, then say \emph{xīnláng} again
\end{itemize}

\begin{tcolorbox}[breakable,colback=teal!4,colframe=teal!35!black,title={Before you move on}]
What does xīnláng mean? (\textquotedblleft{}Bridegroom\textquotedblright{}.) Say it once more.
\end{tcolorbox}
